% =====================================================================================
%  Methods. Formal statement of the predictive layer, the two scoring formulations
%  and the assignment problem. Every constant reported here is fitted, and the file
%  it is fitted in is named.
% =====================================================================================

\section{Methods}
\label{sec:methods}

\begin{figure}[htbp]
\centering
\includegraphics[width=0.98\textwidth]{fig20_system_pipeline.pdf}
\caption{The apparatus end to end. Five event logs enter one code path; the admissibility
filter of Figure~\ref{fig:mechanism} splits each into an unconstrained arm and an admissible
arm, and the gap between what the two arms score is the measurement. The four learners are
held fixed across arms by design, so a difference between arms is a property of the
constraint and not of the model. The scoring formulations and the assignment layer on the
right run on BPI 2014 only, and are what an admissible model is then used for.}
\label{fig:pipeline}
\end{figure}

\subsection{Notation}

Write $i \in \mathcal{I}$ for an incident and $j \in \mathcal{G}$ for an assignment group, with
$n = \lvert \mathcal{I} \rvert$ the number of incidents in an arrival batch and
$m = \lvert \mathcal{G} \rvert$ the number of eligible groups. Case $i$ carries an attribute
vector $\mathbf{x}_i \in \mathbb{R}^{d}$, a binary service level label $y_i$ from
Equation~\ref{eq:label}, an open time $t^{\mathrm{open}}_i$, a first-assignment time
$t^{\mathrm{dec}}_i$ and a resolution time $t^{\mathrm{res}}_i$. The decision variable
$z_{ij} \in \{0,1\}$ is one when incident $i$ is assigned to group $j$. A hat denotes a
quantity estimated from data.

\subsection{Predictive layer}
\label{sec:predictive-methods}

Four gradient-boosted and bagged tree ensembles are evaluated: XGBoost
\citep{chen2016xgboost}, LightGBM \citep{ke2017lightgbm}, CatBoost
\citep{prokhorenkova2018catboost} and a random forest \citep{breiman2001random}. Tree ensembles
rather than deep sequence models are used deliberately.
\citet{fertig2026revisiting} revisit predictive process monitoring under foundation models and
find tabular gradient boosting still competitive, and \citet{grinsztajn2022why} give the
general result that tree ensembles remain ahead of neural networks on tabular data of this
size. Since the object of study here is the effect of a feature-admissibility constraint, the
learner is a control variable and is held fixed across configurations rather than tuned per
configuration.

Class imbalance is handled inside each learner rather than by resampling, so that no synthetic
rows enter a leakage study: XGBoost receives \texttt{scale\_pos\_weight} set to the training
negative-to-positive ratio, LightGBM uses \texttt{is\_unbalance}, CatBoost uses balanced
automatic class weights, and the random forest uses balanced class weights.
Table~\ref{tab:config} gives the full configuration.

\begin{table}[htbp]
\centering
\caption{Experimental configuration. Identical across every configuration and every split, so
that differences between configurations are attributable to the admissibility constraint and
not to tuning.}
\label{tab:config}
\begin{tabular}{lL{9.4cm}}
\toprule
Setting & Value \\
\midrule
XGBoost & 300 trees, max depth 6, learning rate 0.05, subsample 0.8, column subsample 0.8, \texttt{scale\_pos\_weight} $=n_-/n_+$, logloss objective \\
LightGBM & 300 trees, max depth 6, learning rate 0.05, subsample 0.8, column subsample 0.8, \texttt{is\_unbalance} \\
CatBoost & 300 iterations, depth 6, learning rate 0.05, balanced automatic class weights \\
Random forest & 300 trees, max depth 12, balanced class weights \\
\midrule
Primary split & Stratified random, \num{70}/\num{30}, seed 42 \\
Replication split & Stratified random, \num{80}/\num{20}, same seed \\
Causal split & Chronological, earliest \SI{80}{\percent} by open time trains, latest \SI{20}{\percent} evaluates \\
Cross-validation & Stratified 5-fold on the training split, target encoding refitted inside every fold \\
Target-encoding prior & $\alpha = 20$ \\
Decision threshold & 0.5 by default; swept in Section~\ref{sec:threshold}; tuned on a held-out slice of the training period for the chronological arm only \\
Bootstrap resamples & 2000, percentile intervals \\
Seed & 42 throughout; two invocations produce bitwise-identical metrics \\
Hardware & 28-thread x86-64 CPU, no GPU used at any point \\
Software & Python 3.11, scikit-learn, XGBoost, LightGBM, CatBoost, SHAP, pymoo 0.6.2 \\
\bottomrule
\end{tabular}
\end{table}

\subsection{Evaluation metrics}
\label{sec:metrics}

Write TP, FP, TN and FN for the four confusion-matrix cells at a given decision threshold.
The reported quantities are
\begin{align}
\mathrm{Precision} &= \frac{\mathrm{TP}}{\mathrm{TP}+\mathrm{FP}},
&\mathrm{Recall} &= \frac{\mathrm{TP}}{\mathrm{TP}+\mathrm{FN}},
\label{eq:prec}\\[2pt]
\mathrm{Specificity} &= \frac{\mathrm{TN}}{\mathrm{TN}+\mathrm{FP}},
&F_1 &= \frac{2\,\mathrm{TP}}{2\,\mathrm{TP}+\mathrm{FP}+\mathrm{FN}},
\label{eq:f1}
\end{align}
and balanced accuracy
$\mathrm{BA} = \tfrac{1}{2}(\mathrm{Recall}+\mathrm{Specificity})$.
We also report the Matthews correlation coefficient, which is the only one of these that
degrades under \emph{both} kinds of error on \emph{both} classes and is therefore the
single most informative scalar on an imbalanced problem \citep{chicco2020advantages}:
\begin{equation}
\mathrm{MCC} = \frac{\mathrm{TP}\cdot\mathrm{TN} - \mathrm{FP}\cdot\mathrm{FN}}{\sqrt{D}},
\label{eq:mcc}
\end{equation}
where the normalising denominator is
\begin{multline}
D = (\mathrm{TP}+\mathrm{FP})(\mathrm{TP}+\mathrm{FN})\\
\times(\mathrm{TN}+\mathrm{FP})(\mathrm{TN}+\mathrm{FN}).
\label{eq:mccD}
\end{multline}

Threshold-free ranking quality is ROC-AUC, and because the positive class is a
\SI{36.7}{\percent} minority we also report average precision, the area under the
precision-recall curve, whose chance level is the class prevalence rather than \num{0.5}.
Probability quality is the Brier score,
$\mathrm{BS} = N^{-1}\sum_i (\hat{p}_i - y_i)^2$, reported because the scores feed a downstream
decision and a ranking that is correct but miscalibrated still produces the wrong action
\citep{filho2023classifier,perezlebel2025decision}.

Two conventions are used throughout. First, \emph{above-chance signal} means $\mathrm{AUC}-0.5$,
so a drop from \num{0.9050} to \num{0.8258} is a loss of \num{0.0792} out of \num{0.4050}, which
is \SI{19.6}{\percent}. Second, every interval reported is a percentile bootstrap interval over
2000 resamples of the test set, paired across configurations where a difference is quoted.

\subsection{F1: the adaptive assignment score}
\label{sec:f1}

The assignment score $A(i,j)$ maps an incident-group pair to $[0,1]$ by aggregating seven
criteria, each of which is normalised to $[0,1]$ and each of which is either a model output or
a fitted function of observable quantities. Nothing in it is a hand-chosen constant.

\paragraph{$s_1$, outcome propensity.} The isotonic-calibrated probability that the case is
resolved first time right by group $j$, taken from the predictive layer of
Section~\ref{sec:predictive-methods}.

\paragraph{$s_2$, service level survival.} $s_2(i,j) = 1 - \hat{P}(\text{breach} \mid i, j)$,
taken from F2 in Section~\ref{sec:f2}.

\paragraph{$s_3$, expertise match.} Let $\mathbf{r}_i$ be the one-hot requirement vector of the
incident's configuration-item category and $\mathbf{e}_j$ the group's historical specialisation
vector, the distribution of categories it has handled during the training period. Expertise is
the cosine similarity
\begin{equation}
s_3(i,j) \;=\; \frac{\mathbf{e}_j \cdot \mathbf{r}_i}
{\lVert \mathbf{e}_j \rVert_2 \,\lVert \mathbf{r}_i \rVert_2} \;\in\; [0,1].
\label{eq:cosine}
\end{equation}
Cosine rather than a raw count because it is invariant to group volume: a large group that has
handled every category is not thereby a specialist in any of them.

\paragraph{$s_4$, workload headroom.} Over the post-assignment load vector
$\boldsymbol{\ell} = (\ell_1,\dots,\ell_m)$, headroom is Jain's fairness index
\citep{jain1984quantitative}
\begin{equation}
s_4 \;=\; \frac{\left(\sum_{k=1}^{m} \ell_k\right)^2}{m \sum_{k=1}^{m} \ell_k^2}
\;\in\; \left(\tfrac{1}{m},\, 1\right].
\label{eq:jain}
\end{equation}
It is bounded, scale-free and equals one exactly at perfect balance, which the standard
deviation is not and does not.

\paragraph{$s_5$, freshness.} The probability of a favourable outcome decays with the delay
before first touch. Fitting an exponential $\exp(-\lambda t)$ by maximum likelihood on the
training period gives
$\hat{\lambda} = \num{0.0174}\ \text{h}^{-1}$, \SI{95}{\percent} interval
$[\num{0.0161}, \num{0.0188}]$, on $n = \num{31394}$ cases. That is a half-life of
\num{39.9} hours. Cases whose recorded delay exceeds \num{168} hours are held out of this
fit as a different operational regime, which is \num{171} of \num{39449} cases; the uncapped
control fit is reported in Section~\ref{sec:formulation-results}.

\paragraph{$s_6$, urgency.} A logistic ramp in time remaining against the service level target,
$s_6 = \bigl(1 + \exp(k(\tau - \tau_0))\bigr)^{-1}$, with $\hat{k} = \num{0.0228}$ and
$\hat{\tau}_0 = \num{-12.03}$ hours fitted on the same training period.

\paragraph{$s_7$, importance.} A weighted geometric blend of normalised priority, impact and
urgency,
\begin{equation}
s_7(i) \;=\; \prod_{c \in \{\text{pri},\text{imp},\text{urg}\}} \tilde{c}_i^{\,\omega_c},
\qquad \sum_c \omega_c = 1 ,
\label{eq:importance}
\end{equation}
geometric rather than arithmetic so that a low value on one axis cannot be bought back
completely by a high value on another.

\subsubsection{Aggregation by weighted power mean}

The seven criteria are combined by the weighted power mean of order $p$
\citep{dyckhoff1984generalized,dujmovic2015weighted}:
\begin{equation}
A(i,j) \;=\; \Bigl( \sum_{k=1}^{7} w_k \, s_k(i,j)^{\,p} \Bigr)^{1/p},
\qquad \sum_k w_k = 1 .
\label{eq:powermean}
\end{equation}

The exponent is the point of the formulation, because it makes the degree of permitted
compensation between criteria an explicit parameter rather than a hidden assumption. At $p=1$
the mean is arithmetic and fully compensatory, so a strong criterion offsets a weak one; as
$p \to 0$ it converges to the weighted geometric mean; at $p=-1$ it is harmonic; and as
$p \to -\infty$ it converges to the minimum and becomes fully non-compensatory. Reporting a
sensitivity curve over $p$ therefore turns the qualitative claim that one bad criterion should
or should not sink an assignment into a measured one.

\paragraph{Degeneracy, and the floor it forces.} For $p \le 0$ the mean is undefined or
identically zero whenever any criterion is exactly zero. Cosine expertise \emph{is} exactly
zero for a group with no training history in the incident's category, and this occurs in
\num{109012} of \num{348000} decision-matrix cells, which is \SI{31.3}{\percent}. An explicit floor
of $\varepsilon = 10^{-3}$ is therefore applied before any $p \le 0$ aggregation. It is reported
as a parameter of the method rather than buried as an implementation detail, because it changes
the value of the score in a third of all cells.

\subsubsection{Weighting}

Three weighting schemes are computed and all three are reported. CRITIC
\citep{diakoulaki1995determining} derives weight from contrast intensity times conflict. With
the decision matrix min-max normalised to $\mathbf{Z}$, standard deviation $\sigma_k$ per
column and Pearson correlation $\rho_{k\ell}$ between columns,
\begin{equation}
w^{\mathrm{CRITIC}}_k \;\propto\; \sigma_k \sum_{\ell=1}^{7} \bigl(1 - \rho_{k\ell}\bigr).
\label{eq:critic}
\end{equation}
Entropy weighting takes $P_{ik} = Z_{ik} / \sum_i Z_{ik}$, computes
$E_k = -\left(\ln N\right)^{-1}\sum_i P_{ik} \ln P_{ik}$ and sets
$w^{\mathrm{ent}}_k \propto 1 - E_k$. Uniform weighting sets $w_k = 1/7$.

CRITIC alone would be unsafe here, and the reason is structural rather than a matter of taste.
Two of the seven criteria are model outputs, so their spread reflects the calibration of a
classifier rather than the importance of a criterion, and isotonic calibration compresses
variance. Reporting CRITIC alone would launder a modelling artifact into a data-driven weight.
The question is therefore settled empirically, by measuring whether the induced assignment
ordering depends on the scheme at all (Section~\ref{sec:weighting-results}).

\subsection{F2: service level risk as a derived quantity}
\label{sec:f2}

Modelling breach directly as the survival event is degenerate on this log, and this was
established empirically before the design was changed rather than assumed. Because the label is
defined by handle time exceeding a per-priority threshold, a case that has been open longer than
its own threshold has already breached: the measured breach rate past threshold is exactly
\num{1.0000}. A hazard model given elapsed time and priority therefore reads its label rather
than predicting it. An early version of this model scored \num{0.9418} AUC while learning
nothing, and that number is reported here as a worked example of the failure mode the paper is
about.

The admissible object is the hazard of \emph{resolution}, which elapsed time does not determine.
Each case contributes one row per elapsed-time bin it survives into, over the ten bins with
edges at $0, 1, 2, 4, 8, 16, 24, 48, 96, 168$ hours and an open final bin. The expansion turns
\num{39449} cases into \num{161720} person-period rows, a factor of \num{4.10}, with a per-bin
resolution rate of \num{0.2430}. A gradient-boosted classifier estimates the discrete hazard
\begin{equation}
h(t \mid \mathbf{x}) \;=\; P\bigl(t^{\mathrm{res}} \in b_t \;\big|\; t^{\mathrm{res}} \ge b_t,\, \mathbf{x}\bigr)
\label{eq:hazard}
\end{equation}
directly, which is the discrete-time survival formulation of \citet{suresh2022survival}.
Survival follows by the product rule and breach risk is the conditional survival ratio:
\begin{align}
S(t \mid \mathbf{x}) &= \prod_{u \le t} \bigl(1 - h(u \mid \mathbf{x})\bigr),
\label{eq:surv}\\[2pt]
P\bigl(\text{breach} \mid \text{alive at } t, \mathbf{x}\bigr)
  &= \frac{S(\tau \mid \mathbf{x})}{S(t \mid \mathbf{x})},
\label{eq:survival}
\end{align}
where $\tau$ is the case's own service level threshold from Equation~\ref{eq:label}. Monotonicity
of $S$ is asserted as a test rather than assumed, and the test passes.

The covariates are priority, impact, urgency, assignment delay, category, the bin index and the
bin start time, all of which satisfy Equation~\ref{eq:rule1}. The escalation cut point is tuned
on training-period data alone and then frozen at \num{0.5305}; tuning it on the evaluation
period would reproduce, on the time axis, exactly the leak this paper measures on the feature
axis.

\subsection{Multi-objective assignment}
\label{sec:moo}

A batch of $n$ incidents must be assigned across $m$ eligible groups. The formulation is
\begin{align}
\max\; & f_1(\mathbf{z}) = \tfrac{1}{n}\textstyle\sum_{i,j} z_{ij}\, \hat{p}^{\mathrm{ftr}}_{ij}
  \label{eq:f1obj}\\
\min\; & f_2(\mathbf{z}) = \tfrac{1}{n}\textstyle\sum_{i,j} z_{ij}\, \hat{p}^{\mathrm{breach}}_{ij}
  \label{eq:f2obj}\\
\min\; & f_3(\mathbf{z}) = \operatorname{sd}\bigl(\ell_1,\dots,\ell_m\bigr),
  \;\; \ell_j = \textstyle\sum_i z_{ij} \label{eq:f3obj}\\
\max\; & f_4(\mathbf{z}) = \tfrac{1}{n}\textstyle\sum_{i,j} z_{ij}\, s_3(i,j)
  \label{eq:f4obj}\\
\min\; & f_5(\mathbf{z}) = \tfrac{1}{n}\textstyle\sum_{i,j} z_{ij}\, \hat{d}_{ij}
  \label{eq:f5obj}
\end{align}
where $f_1$ is expected first-time-right resolution, $f_2$ expected service level breach, $f_3$
workload imbalance, $f_4$ expertise match and $f_5$ expected delay before first touch. The
formulation is subject to
\begin{align}
\text{C1:}\;\; & \textstyle\sum_{j} z_{ij} = 1 && \forall i \label{eq:c1}\\
\text{C2:}\;\; & \textstyle\sum_{i} z_{ij} \le \kappa_j && \forall j \label{eq:c2}\\
\text{C3:}\;\; & z_{ij} \le a_j && \forall i,j \label{eq:c3}\\
\text{C4:}\;\; & z_{ij} \, s_3(i,j) \ge \eta \, z_{ij} && \forall i,j \label{eq:c4}
\end{align}
with $z_{ij}\in\{0,1\}$. C1 assigns each incident exactly once, C2 respects the capacity
$\kappa_j$ of group $j$, C3 restricts assignment to available groups through the indicator
$a_j$, and C4 imposes a minimum expertise eligibility $\eta$. Constraint C2 is enforced by a bounded repair operator rather than by
penalty, so that no infeasible solution can enter the population; the operator is capped at six
passes and its effect is verified by counting violations, which are zero across all
\num{1225} runs.

Workload imbalance is kept as the standard deviation in $f_3$. Jain's index of
Equation~\ref{eq:jain} is reported \emph{alongside} it for every policy and never substituted
for it, because changing an objective after observing that a method loses on it would be
fitting the metric to the result.

\subsubsection{The proposed operators}

Three modifications to NSGA-II \citep{deb2002fast} are proposed and, critically, are ablated
individually rather than reported as a bundle.

\begin{description}
\item[Priority-aware initialisation.] A fraction of the initial population is constructed
greedily in descending order of incident priority rather than sampled uniformly, so that the
search starts from solutions in which urgent cases already sit with capable groups.
\item[Adaptive crossover.] The crossover probability rises when population diversity, measured
as the mean pairwise assignment Hamming distance, falls below a threshold.
\item[Service-level-aware mutation.] Mutation is biased towards reassigning the incidents with
the highest current breach probability.
\end{description}

\begin{algorithm}[htbp]
\caption{Priority-aware adaptive NSGA-II (PAA-NSGA-II). Differences from standard NSGA-II are
lines 1, 6 and 7.}
\label{alg:paa}
\begin{algorithmic}[1]
\State $P_0 \gets$ \textsc{PriorityAwareInit}$(N, \rho)$ \Comment{$\rho$ of $N$ built greedily by descending priority}
\State $P_0 \gets$ \textsc{CapacityRepair}$(P_0)$
\For{$g = 1$ \textbf{to} $G$}
  \State $\delta \gets$ \textsc{Diversity}$(P_{g-1})$ \Comment{mean pairwise Hamming distance}
  \State $Q \gets$ \textsc{TournamentSelect}$(P_{g-1})$
  \State $Q \gets$ \textsc{Crossover}$(Q, \; p_c(\delta))$ \Comment{$p_c$ raised when $\delta$ is low}
  \State $Q \gets$ \textsc{Mutate}$(Q, \; p_m,\; \text{bias} \propto \hat{p}^{\mathrm{breach}})$
  \State $Q \gets$ \textsc{CapacityRepair}$(Q)$
  \State $R \gets P_{g-1} \cup Q$
  \State $\mathcal{F} \gets$ \textsc{FastNonDominatedSort}$(R)$ \Comment{$O(MN^2)$}
  \State $P_g \gets$ \textsc{SelectByRankAndCrowding}$(\mathcal{F}, N)$
\EndFor
\State \Return non-dominated front of $P_G$
\end{algorithmic}
\end{algorithm}

Hypervolume is computed against a reference point fixed as the nadir across all variants and all
runs within a batch, in a first pass completed before any hypervolume is taken. This matters:
the alternative, taking the reference point from whichever run finished first, makes every
hypervolume value depend on execution order, and correcting it reversed the sign of this study's
headline comparison (Section~\ref{sec:harness-defects}).
